% Part: normal-modal-logic
% Chapter: filtrations
% Section: introduction

\documentclass[../../../include/open-logic-section]{subfiles}

\begin{document}

\olfileid[hi]{nml}{fil}{int}

\olsection{परिचय}

किसी तर्कशास्त्र के बारे में एक महत्त्वपूर्ण प्रश्न सदा यह होता है कि क्या
वह निर्णेय है, अर्थात् क्या ऐसी प्रभावी प्रक्रिया है जो इस प्रश्न का उत्तर
देगी: ``क्या यह !!{formula} मान्य है?'' प्रतिज्ञप्तिक तर्कशास्त्र निर्णेय है:
हम सत्यता-सारणी बनाकर प्रभावी ढंग से जाँच सकते हैं कि कोई !!a{formula}
सर्वसत्य है या नहीं, और किसी दिए गए !!{formula} की सत्यता-सारणी परिमित होती
है। किंतु हम सीधे यह जाँच नहीं सकते कि कोई मोडल सूत्र सभी प्रतिरूपों में
सत्य है या नहीं, क्योंकि प्रतिरूप अपरिमिततः अनेक हैं। हम किसी दिए गए सूत्र
से संबद्ध सभी परिमित प्रतिरूपों को सूचीबद्ध कर सकते हैं, क्योंकि केवल उन
!!{propositional variable}ों को जगतों के उपसमुच्चय नियत करना प्रासंगिक है जो
वास्तव में !!{formula} में आते हैं। यदि अभिगम्यता संबंध नियत है, तो संभावित
भिन्न नियतन $V(p)$ केवल $W$ के सभी उपसमुच्चय हैं; और यदि $\card{W} = n$,
तो ऐसे $2^n$ नियतन हैं। यदि हमारे !!{formula}~$!A$ में $m$ !!{propositional variable} हैं, तो $n$ जगतों वाले $2^{nm}$ भिन्न प्रतिरूप हैं। प्रत्येक प्रतिरूप
के लिए हम हर जगत पर $!A$ का सत्यता-मूल्य गणित करके जाँच सकते हैं कि $!A$
सभी जगतों पर सत्य है या नहीं। निस्संदेह हमें सभी संभावित अभिगम्यता संबंधों
को भी जाँचना है, किंतु $n$ जगतों पर भी परिमिततः अनेक संबंध ही होते हैं
(विशेषतः $W \times W$ के उपसमुच्चयों की संख्या, अर्थात् $2^{n^2}$।

यदि हमारी रुचि तर्कशास्त्र~\Log{K} में नहीं, बल्कि प्रतिरूपों के किसी वर्ग
(जैसे स्वपरावर्ती संक्रमणशील प्रतिरूपों) द्वारा परिभाषित तर्कशास्त्र में
है, तो हमें यह भी जाँच सकना चाहिए कि अभिगम्यता संबंध सही प्रकार का है या
नहीं। जब भी हमारी रुचि के फ़्रेम मोडल !!{formula}ों द्वारा परिभाष्य हों, हम
ऐसा कर सकते हैं (जैसे यह जाँचकर कि \Ax{T} और \Ax{4} फ़्रेम में मान्य हैं
या नहीं)। अतः विचार यह होगा कि सभी परिमित फ़्रेमों को क्रमशः लें, हर एक के
लिए जाँचें कि वह हमारी रुचि के वर्ग का फ़्रेम है या नहीं, फिर उस फ़्रेम पर
सभी संभावित प्रतिरूप सूचीबद्ध करें और जाँचें कि प्रत्येक में $!A$ सत्य है
या नहीं। यदि नहीं, तो रुकें: $!A$ रुचिकर प्रतिरूप-वर्ग में मान्य नहीं है।

इस विचार में एक समस्या है: हम नहीं जानते कि हम खोजना कब रोक सकते हैं, यदि
कभी रोक भी सकें। यदि !!{formula} का कोई परिमित प्रतिप्रतिरूप है, तो हमारी
प्रक्रिया उसे खोज लेगी। किंतु यदि उसका कोई परिमित प्रतिप्रतिरूप नहीं है,
तो हमें उत्तर नहीं मिलेगा। !!{formula} मान्य हो सकता है (कोई प्रतिप्रतिरूप
ही न हो), या उसका केवल अपरिमित प्रतिप्रतिरूप हो सकता है, जिसे हम कभी
देखेंगे ही नहीं। यदि हम दिखा सकें कि प्रतिप्रतिरूप वाले हर !!{formula} का
कोई परिमित प्रतिप्रतिरूप भी होता है, तो इस समस्या को दूर किया जा सकता है।
ऐसी स्थिति में हम कहते हैं कि तर्कशास्त्र में \emph{परिमित प्रतिरूप गुण}
है।

पर हम कैसे दिखाएँ कि किसी तर्कशास्त्र में परिमित प्रतिरूप गुण है? इसका एक
उपाय $!A$ के किसी अपरिमित (प्रति)प्रतिरूप को परिमित प्रतिरूप में बदलने की
विधि खोजना होगा। यदि ऐसा हो सके, तो जब भी कोई ऐसा प्रतिरूप हो जिसमें $!A$
सत्य नहीं है, उससे मिला परिमित प्रतिरूप भी $!A$ को असत्य बनाएगा। वह
परिमित प्रतिरूप हमारी सभी परिमित प्रतिरूपों की सूची में आएगा और अंततः हम
हर अमान्य सूत्र के लिए निर्धारित कर लेंगे कि वह अमान्य है। यदि सूत्र मान्य
है, तो हमारी प्रक्रिया समाप्त नहीं होगी। यदि इसके अतिरिक्त हम दिखा सकें
कि हमारी प्रक्रिया से मिलने वाले परिमित प्रतिरूप के आकार की कोई अधिकतम
सीमा है, और यह अधिकतम आकार केवल !!{formula}~$!A$ पर निर्भर करता है, तो हमें
वह आकार मिल जाएगा जहाँ तक प्रतिप्रतिरूप खोजते हुए परिमित प्रतिरूपों को
जाँचना है। यदि तब तक प्रतिप्रतिरूप न मिला, तो कोई है ही नहीं। तब हमारी
प्रक्रिया वास्तव में हर !!{formula}~$!A$ के लिए ``क्या $!A$ मान्य है?''
प्रश्न को निर्णीत करेगी।

अपरिमित संरचनाओं को परिमित संरचनाओं में बदलने के लिए प्रायः सफल होने वाली
एक रणनीति संरचना के उन !!{element}ों की ``पहचान'' करना है जो प्रासंगिक
पहलुओं में समान व्यवहार करते हैं। यदि $\mModel{M}$ में अपरिमिततः अनेक जगत
प्रासंगिक पहलुओं में समान व्यवहार करते हैं, तो संभव है कि हम \emph{सभी}
जगतों को ऐसे जगतों के परिमिततः अनेक (संभवतः अपरिमित) ``वर्गों'' में समेट
सकें। दूसरे शब्दों में, हमें जगतों के समुच्चय का उचित ढंग से विभाजन करना
चाहिए: ऐसा कि हर विभाजन में अपरिमिततः अनेक जगत हों, किंतु विभाजन केवल
परिमिततः अनेक हों। फिर हम नया प्रतिरूप~$\mModel{M^*}$ परिभाषित करते हैं
जिसमें ये विभाजन ही जगत हैं। पुराने प्रतिरूप के परिमिततः अनेक विभाजन हमें
नए प्रतिरूप में परिमिततः अनेक जगत देते हैं, अर्थात् परिमित प्रतिरूप। जिस
विभाजन में जगत~$w$ है, उसे $[w]$ कहें। हम चाहेंगे कि
$\mSat{M}{!A}[w]$ तभी जब $\mSat{M^*}{!A}[{[w]}]$, क्योंकि यदि पुराना
प्रतिरूप $!A$ का प्रतिप्रतिरूप था, तो हम नए प्रतिरूप को भी उसका
प्रतिप्रतिरूप चाहते हैं। इसके लिए विभाजन और $\mModel{M^*}$ के अभिगम्यता
संबंध, दोनों को उचित ढंग से परिभाषित करना होगा।

यह कैसे चलेगा, इसे देखने के लिए पहले कल्पना करें कि हमारे पास कोई
अभिगम्यता संबंध नहीं है। $\mSat{M}{\Box !B}[w]$ तभी जब किसी $v \in W$
के लिए $\mSat{M}{\Box !B}[v]$; और $\mModel{M^*}$ के लिए भी यही, बस $w$
और $v$ के स्थान पर $[w]$ और $[v]$। पहले विचार के रूप में कहें कि दो जगत
$u$ और $v$ तुल्य हैं (एक ही विभाजन में हैं) यदि वे $\mModel{M}$ के सभी
!!{propositional variable}ों पर सहमत हों, अर्थात् $\mSat{M}{p}[u]$ तभी जब
$\mSat{M}{p}[v]$। $V^*(p) = \Setabs{[w]}{\mSat{M}{p}[w]}$ रखें। हमारा
लक्ष्य दिखाना है कि $\mSat{M}{!A}[w]$ तभी जब
$\mSat{M^*}{!A}[{[w]}]$। स्पष्टतः हम इसे आगमन से सिद्ध करेंगे: आधार स्थिति
$!A \equiv p$ होगी। पहले मान लें $\mSat{M}{p}[w]$। तब परिभाषा से $[w]
\in V^*$, अतः $\mSat{M^*}{p}[{[w]}]$। अब मान लें कि
$\mSat{M^*}{p}[{[w]}]$। इसका अर्थ है कि $[w] \in V^*(p)$, अर्थात् $w$
के तुल्य किसी $v$ के लिए $\mSat{M}{p}[v]$। किंतु ``$w$, $v$ के तुल्य
है'' का अर्थ है ``$w$ और $v$ ठीक उन्हीं !!{propositional variable}ों को सत्य
बनाते हैं,'' अतः $\mSat{M}{p}[w]$। अब आगमनात्मक चरण के लिए, जैसे $!A
\ident \lnot !B$। तब $\mSat{M}{\lnot !B}[w]$ तभी जब
$\mSat/{M}{!B}[w]$, तभी जब $\mSat/{M^*}{!B}[{[w]}]$ (आगमनात्मक
परिकल्पना से), तभी जब $\mSat{M^*}{\lnot !B}[{[w]}]$। अन्य गैर-मोडल
संकारकों के लिए भी इसी प्रकार। यह $\Box$ के लिए भी चलता है: मान लें
$\mSat{M^*}{\Box !B}[{[w]}]$। इसका अर्थ है कि प्रत्येक $[u]$ के लिए
$\mSat{M^*}{!B}[{[u]}]$। आगमनात्मक परिकल्पना से प्रत्येक $u$ के लिए
$\mSat{M}{!B}[u]$। फलतः $\mSat{M}{\Box !B}[w]$।

सामान्य स्थिति में, जहाँ हमें $\mModel{M^*}$ का अभिगम्यता संबंध भी
परिभाषित करना है, बात अधिक जटिल है। हम किसी प्रतिरूप~$\mModel{M^*}$ को
\emph{निस्यंदन} कहेंगे यदि उसका अभिगम्यता संबंध~$R^*$ ऊपर का आगमनात्मक
प्रमाण चलाने के लिए आवश्यक शर्तें संतुष्ट करता हो। तब कोई भी
निस्यंदन~$\mModel{M^*}$, $[w]$ पर $!A$ को तभी सत्य बनाएगा जब
$\mModel{M}$, $w$ पर $!A$ को सत्य बनाता है। किंतु अब हमें यह भी दिखाना
है कि निस्यंदन \emph{होते हैं}, अर्थात् हम $R^*$ को इस प्रकार परिभाषित
कर सकते हैं कि वह आवश्यक शर्तें संतुष्ट करे। इसके चलने के लिए हमें यह
अपेक्षा करनी होगी कि जगत $u$, $v$ केवल सभी !!{propositional variable}ों पर
सहमत होने पर ही नहीं, बल्कि $!A$ के सभी उप-!!{formula}ों पर सहमत होने पर
तुल्य गिने जाएँ। चूँकि $!A$ के केवल परिमिततः अनेक उप-!!{formula} हैं, इससे
अब भी प्रत्याभूत होगा कि निस्यंदन परिमित है। निस्यंदन परिभाषित करने का
केवल एक ही ढंग नहीं है; और यह सुनिश्चित करने के लिए कि निस्यंदन का
अभिगम्यता संबंध अपेक्षित गुण (जैसे स्वपरावर्तिता, संक्रमणशीलता, इत्यादि)
संतुष्ट करे, हमें $R^*$ की परिभाषा में कल्पनाशील होना होगा।


\end{document}


